% Paper CSR — Abstract & Structure
% Target: IEEE Aerospace 2027 (deadline ~Nov 2026) or DASC 2027

\documentclass[conference]{IEEEtran}
\usepackage{cite,amsmath,amssymb,graphicx,hyperref}

\title{Counterfactual Safety Reasoning: Post-Mission What-If Proofs for Autonomous Drones}
\author{Andre Oliveira\\RNFC Project\\\texttt{handre.oliveira@gmail.com}}

\begin{document}
\maketitle

\begin{abstract}
We present Counterfactual Safety Reasoning (CSR), a technique that
enables a drone to answer post-mission ``what if?' questions with
cryptographically verifiable answers. Given a recorded trajectory and
a hash chain of safety events, CSR replays the mission under modified
conditions (e.g., wind $2\times$, altitude $+150$~m) and produces a
proof bound to the original chain. A regulator can verify that
``if wind had been $2\times$, the drone would have missed waypoint~3
by 40~m' without trusting the operator. We demonstrate CSR on
a 4-scenario case study (wind $0\times$, $1\times$, $20\times$,
altitude $+150$~m), showing 4/4 proofs verified and 3/3 attacks
(tampered chain, tampered result, wrong key) rejected. CSR is
implemented in Rust and evaluated on 12 unit tests. We discuss
how CSR complements static safety cases and its limitations
(synthetic physics, FNV-1a hash, no formal guarantees).
\end{abstract}

\begin{IEEEkeywords}
counterfactual reasoning, safety certification, hash chain, UAV,
post-mission analysis, audit
\end{IEEEkeywords}

\section{Introduction}
% 1. Problem: safety cases say "the drone was safe", not "why"
% 2. Gap: no way to prove what would have happened under different conditions
% 3. Contribution: CSR generates verified counterfactual proofs
% 4. Results: 4/4 proofs, 3/3 attacks rejected

\section{Background}
\subsection{Safety Cases and Black-Box Analysis}
% Static safety cases, black-box logs (ArduPilot, PX4)
\subsection{Cryptographic Audit}
% Hash chains, Merkle trees, transparency logs
\subsection{Counterfactual Reasoning}
% Causal inference, sensitivity analysis (offline)

\section{CSR Design}
\subsection{Trajectory Replay}
% Re-run physics with modified conditions
\subsection{Counterfactual Engine}
% what_if_wind(scale), what_if_altitude(offset)
\subsection{Proof Construction}
% proof = SHA-256(chain_root || result_hash || scenario)

\section{Implementation}
\subsection{Crate Structure}
% rnfc-counterfactual (250 lines), 12 tests
\subsection{Integration with CertChain}
% Original chain + counterfactual proofs stored together

\section{Evaluation}
\subsection{Case Study: 4 Scenarios}
% wind 0x (Better), 1x (Same), 20x (Failed), altitude+150m (Failed)
\subsection{Security Properties}
% Tamper chain -> BadSignature; tamper result -> BadSignature
% wrong key -> BadSignature
\subsection{Limitations}
% Synthetic physics; FNV-1a hash (not SHA-256); no formal guarantees

\section{Discussion}
\subsection{Relation to DO-326A}
% CSR provides evidence for post-mission audit (SEC-12)
\subsection{Ethical and Legal Considerations}
% Who owns the counterfactual? Who can request it?

\section{Conclusion}
% CSR enables verified post-mission what-if analysis. Future:
% real physics, formal verification, standardisation.

\bibliographystyle{IEEEtran}
\bibliography{references}

\end{document}
